\documentclass[12pt]{beamer}
\usepackage{beamerthemeCopenhagen}
\usecolortheme[rgb={0.18,0.3,0.3}]{structure}
%\usepackage[T1]{fontenc}
\usepackage[utf8]{inputenc}
%\usepackage[latin1]{inputenc}
%\usepackage[brazil]{babel}
\usepackage{amsmath}
\usepackage{amsfonts}
\usepackage{amsthm}
\usepackage{pgf,tikz}
\usepackage{bbm}

\usepackage{graphicx}
\usepackage{url,color,ae}
\usepackage{listings,color,upquote}
\usepackage{verbatim}
\usepackage{multirow}
\usepackage{booktabs}

\usetikzlibrary{arrows}

\newcommand{\caixa}{$\Box$}
\newtheorem{defi}{Definição}
\newtheorem{ex}{Exemplo}
\newtheorem{teo}{Teorema}
\newtheorem{cor}{Corolário}
\newtheorem{lema}{Lema}
\newtheorem{dist}{Distribuição}
\newtheorem{emv}{Máxima verossimilhança}
\newtheorem{mo1}{Primeiro momento}
\newtheorem{mo2}{Segundo momento}
\newtheorem{prova}{Demonstração}
\newtheorem{BiNeg}{Definição}
\setbeamertemplate{navigation symbols}{}

\institute[IC-UFAL]{Universidade Federal de Alagoas\\
Instituto de Computação}
\title{Comparação de estimadores pelo método Monte Carlo das distribuições: $\mathcal{G}_{\rm I}^ {0}$, Uniforme, Binomial negativa, Gamma e Exponencial}
\author{Antônio Marcos Larangeiras\and Alisson Nascimento\and
  Paulo\and Bruno Normande\and Juliana Leal}
\date{{\tiny Maceió-AL, Dezembro de 2013}}
\begin{document}
\begin{frame}
\maketitle
\begin{picture}(0,0)(0,-50)
\includegraphics[scale=0.11]{IC.png}
\end{picture}
\begin{picture}(0,0)(-250,-40)
\includegraphics[scale=0.13]{ufal.png}
\end{picture}
\end{frame} 

%=========================
\section{Introdução}
%=========================


%================================================================================
\begin{frame}
\frametitle{Introdução}
\begin{itemize}
\item Calcular o estimador por máxima verossimilhança;

\bigskip

\item Calcular estimador pelo primeiro momento;

\bigskip

\item Calcular estimador pelo segundo momento;

\bigskip

\item Utilizar o Método Monte Carlo força bruta.

\end{itemize}
\end{frame}
%===============================================================================


%========================
\section{Distribuições}
%========================

\subsection{Distribuição Uniforme}
%========================================================================================================================================================
\begin{frame}
\frametitle{Distribuição Uniforme $U(0, \theta)$}
\begin{dist}[$U(0, \theta)$]
Denotaremos esta distribuição por $X\sim U(x; 0,\theta)$, com $\theta>0$. Sua densidade é

{\small $$f(x) =
\left \{
\begin{array}{cc}
\frac{1}{ \theta } & 0 \leqslant x \leqslant \theta\\
0 & \textnormal{caso contrário.} \\
\end{array}
\right.$$ }

\bigskip\bigskip

{\scriptsize $E(X)=\dfrac{\theta}{2}$\hspace{1cm} e \hspace{1cm}$Var(X)=\dfrac{\theta^{2}}{12}$ }
\end{dist}
\end{frame}
%========================================================================================================================================================

%====================
\subsubsection{Estimadores}
%====================

\begin{frame}
\frametitle{Estimadores da distribuiç\~ao $U(0, \theta)$}
\begin{emv}[$U(0, \theta)$]
$$\widehat{\theta}_{1}={\dfrac{1}{\theta ^ n }}{\mathbb{1}_{\max{\{X_1,\dots,X_n\}\leqslant\theta}}},
\textnormal{portanto } \widehat{\theta} ={\max{\{X_1,\dots,X_n\}\leqslant\theta}} $$
\end{emv}

\begin{mo1}[$U(0, \theta)$]
$$\widehat{\theta}_{2}=2\overline{X}$$
\end{mo1} 

\end{frame}

\begin{frame}
\begin{mo2}[$U(0, \theta)$]
$$\widehat{\theta}_{3} = \sqrt{{{{3} \over {n}} {\sum_{i=1}^{n} X_{i}^{2}}}} $$
\end{mo2}
\end{frame}
%========================================================================================================================================================
%===================
\subsubsection{Gr\' aficos}
%===================
%=========================================================================================================================================================
\begin{frame}
\frametitle{Gráficos da distribuiç\~ao $U(0, \theta)$}

\end{frame}
%=========================================================================================================================================================


%==============================================================================================
\subsection{Distribuição $\mathcal{G}_{\rm I}^{0}$}

\begin{frame}
\frametitle{Distribui\c c\~ ao $\mathcal{G}_{\rm I}^{0}$}

Denotaremos esta distribuição por $Z\sim\mathcal{G}_{\rm I}^{0}(z; \alpha,1,1)$, com $-\alpha, z>0$. Sua densidade é

\begin{displaymath}
f_{Z}=(z;\alpha,1,1)=\dfrac{\Gamma(L-\alpha)}{\Gamma(-\alpha)(1+z)^{1-\alpha}}, \mbox{com} -\alpha \ \mbox{e}\ z>0.
\end{displaymath}


\end{frame}

%========================================================================================================================================================
\begin{frame}
\frametitle{Distribui\c c\~ ao $\mathcal{G}_{\rm I}^{0}$}

A esperança matemática da distribuição $\mathcal{G}_{\rm I}^{0}$:

\begin{equation}
\label{eq:Gi0_esp}
E(Z)=\dfrac{\gamma}{\alpha-1}
\end{equation}\pause

A variância da distribuição $\mathcal{G}_{\rm I}^{0}$:

\begin{equation}
\label{eq:Gi0_var}
Var(Z)=\dfrac{1}{(\alpha-1)^2}+\dfrac{1}{(\alpha-2)(\alpha-1)^2}+\dfrac{1}{\alpha-1}
\end{equation}

\end{frame}

%====================
\subsubsection{Estimadores}
%====================
\begin{frame}
\frametitle{Estimadores da Distribuição $\mathcal{G}_{\rm I}^{0}$}
Pela Máxima Veressimilhança:
\begin{equation}
\label{eq:Gi0_MV}
\hat{\alpha}=\dfrac{n}{\displaystyle\sum_{i=1}^{n}\ln(1+z_i)}
\end{equation} \pause
Pelo momento central de ordem 2:
\begin{eqnarray}
\label{eq:Gi0_m_p2}
\hat{\alpha}_2=\overline{Z}
\end{eqnarray}
\end{frame}

\begin{frame}
\frametitle{Estimadores da Distribuição $\mathcal{G}_{\rm I}^{0}$}
Pelo segundo momento:
\begin{equation}
\label{eq:Gi0_m_p3}
\hat{\alpha}_3 = \frac{\frac{3}{k}\displaystyle\sum_{i=1}^{k}Z_{i}^{2}+2+\sqrt{\frac{1}{k^{2}}\displaystyle\sum_{i=1}^{k}Z_{i}^{4}+4}}{\frac{2}{k}\displaystyle\sum_{i=1}^{k}Z_{i}^{2}}
\end{equation}

\end{frame}


%===================
\subsubsection{Gráficos}
%=========================================================================================================================================================
\begin{frame}
\frametitle{Gráficos da distribuição $\mathcal{G}_{\rm I}^{0}$}

\end{frame}
%=========================================================================================================================================================




\subsection{Distribuição Binomial Negativa}
\begin{frame}
\frametitle{Distribuição Binomial Negativa}
Considere a situação de observar um evento dicotômico $Y_{i}$ independentes e identicamente distribuídos segundo uma lei de Bernoulli de probabilidade $p$. Suponha que se registre $Y$, o número de ensaios até obter exatamente $k$ sucessos. 

\end{frame}

\begin{frame}
\frametitle{Distribuição Binomial Negativa}

Seja uma variável aleatória que fornece o numero de ensaios até o k-ésimo sucesso. 
Assim, $Y$ tem uma distribuição binomial negativa com parâmetro $p \in (0,1)$, 
se sua função de probabilidade é dada por:

\begin{displaymath}
P_{r}(Y=y)=\left\{ \begin{array}{ll}
\binom{y-1}{k-1} \cdot p^{k} \cdot (1-p)^{y-k} & \textrm{se  y=k,k+1, \ldots.}\\
0 & \textrm{, caso contrário}
\end{array} \right.
\end{displaymath}


\end{frame}

%========================================================================================================================================================
\begin{frame}
\frametitle{Distribui\c c\~ ao $BN(p,k)$}

Usualmente sua função de probabilidade denota-se: $Y \sim BN(p,k)$.\pause

A esperança matemática da distribuição binomial negativa:

\begin{equation}
\label{eq:esp}
E(X) = {\frac{k}{p}}
\end{equation}\pause

A variância da distribuição binomail negativa:

\begin{equation}
\label{eq:var}
Var(X) = {\frac{k(1-p)}{p^{2}}}
\end{equation}

\end{frame}

%====================
\subsubsection{Estimadores}
%====================
\begin{frame}
\frametitle{Estimadores da Distribuição $BN(p,k)$}
Pela Máxima Veressimilhança:
\begin{equation}
\label{eq:BN_MV}
{\hat{p}} = {\frac{n k}{\sum_{i=1}^{n}{Y_{i}}}}
\end{equation} \pause
Pelo momento central de ordem 2:
\begin{eqnarray}
\label{eq:BN_m_p2}
{\hat{p}_{2}} = {\frac{-k + \sqrt{k^{2}+4k{Var(X_{i})}}}{2{Var(X_{i})}}}
\end{eqnarray}
\end{frame}

\begin{frame}
\frametitle{Estimadores da Distribuição $BN(p,k)$}
Pelo segundo momento:
\begin{equation}
\label{eq:BN_m_p3}
{\widehat{p}_{3}} = {\frac{{-k} + {\sqrt{k^{2} + 4 E(X_{i}^{2})(k^{2}+k)}}}{2E(X_{i}^{2})}}
\end{equation}

\end{frame}


%===================
\subsubsection{Gráficos}
%=========================================================================================================================================================
\begin{frame}
\frametitle{Gráficos da distribuição $BN(p,k)$}

\end{frame}
%=========================================================================================================================================================


\subsection{Distribuição Gamma}

% =====================================================================================
 \begin{frame}
   \frametitle{Distribuição Gamma}

   A Distribuição Gamma é caracterizada por dois valores, denominados
   \textit{shape} ($k$) e \textit{scale} ($\theta$).

   \begin{dist}[Função Probabilidade Densidade]
     $$ f(x;k,\theta) = \frac{x^{k-1}e^{-\frac{x}{\theta}}}{\theta^k\Gamma(x)} $$ 

     \begin{center}
       para   

       $ x > 0 $ e $ k,\theta > 0 $ 
     \end{center}
   \end{dist}
 \end{frame}
% ======================================================================================

% ====================
\subsubsection{Estimadores}
% ====================


\begin{frame}
  \frametitle{Estimadores da distribuição Gamma}

  Como foi definido que o valor de $\theta$ seria sempre igual a 1,
  então basta estimar os valores de $k$.

  \begin{emv}[Gamma]
    
  \end{emv}
\end{frame}
% ===================================================================================
\begin{frame}
  \frametitle{Estimadores da distribuição Gamma}
  \begin{mo1}[Gamma]
    \begin{align*}
      E[X] &= k\theta \\
      \hat{k} &= \frac{1}{n} \sum_{i=1}^nx_i
    \end{align*}
  \end{mo1}
  
  \begin{mc2}[Gamma]
    \begin{align*}
      Var[X] &= k\theta^2 \\
      \hat{k} &= Var(x)
    \end{align*}
  \end{mc2}
\end{frame}
% ========================================================================================================================================================

% ===================
\subsubsection{Gráficos}
% ===================
% =========================================================================================================================================================
\begin{frame}
  \frametitle{Gráficos da distribuição $\mathcal{G}_{\rm I}^{0}$}

\end{frame}
% =========================================================================================================================================================

%==============================================================================================
\subsection{Distribuição Exponencial}
\begin{frame}
\frametitle{Distribuição Exponencial}
A distribuição exponencial é frequentemente usada em estudos de confiabilidade como sendo um modelo para o tempo até a falha de um equipamento. Essa distribuição funciona como o inverso da distribuição Poisson. Enquanto a Poisson estima a quantidade de eventos em um intervalo, a exponencial analisa um intervalo ou espaço para a ocorrência de um evento. 

Sua função densidade é dada por:

\begin{equation}
 F(v) = 1 - e^{-\lambda v}, \textnormal{para } v \geq 0.
\end{equation}

\end{frame}

%========================================================================================================================================================
\begin{frame}
\frametitle{Distribuição $Exp(\lambda)$}

A esperança matemática da distribuição $Exp(\lambda)$:

\begin{equation}
\label{eq:Exp_esp}
E(V)=1/\lambda
\end{equation}\pause

A variância da distribuição $Exp(\lambda)$:

\begin{equation}
\label{eq:Esp_var}
Var(V)=1/\lambda^2
\end{equation}

\end{frame}

%====================
\subsubsection{Estimadores}
%====================
\begin{frame}
\frametitle{Estimadores da Distribuição $Exp(\lambda)$}
Pela Máxima Veressimilhança:
\small \begin{equation}
\label{eq:Exp_MV}
\hat{\lambda}=1/\bar{V}
\end{equation} \pause
Pelo momento primeiro momento:
\begin{eqnarray}
\label{eq:Exp_m_p2}
\hat{\lambda}_{2}=1/\lambda
\end{eqnarray}
\end{frame}

\begin{frame}
\frametitle{Estimadores da Distribuição $Exp(\lambda)$}
Pelo segundo momento:
\begin{equation}
\label{eq:Exp_m_p3}
\hat{\lambda}_{3} = 2/\lambda^2
\end{equation}

\end{frame}


%===================
\subsubsection{Gráficos}
%=========================================================================================================================================================
\begin{frame}
\frametitle{Gráficos da distribuição $Exp(\lambda)$}

\end{frame}
%=========================================================================================================================================================





%===================
\section{Resultados}
%===================

%================================================================================
%\input{../unforme/result_uniforme}
\subsection{Distribuição Binomial Negaiva}
\begin{frame}
\frametitle{Distribuição $BN(p,k)$}
\tiny
\begin{table}[h]
\caption{Viés dos estimadores $\hat{p_{1}}$ e $\tilde{p_{1}}$.}
\label{tab:p1Vies}
\centering
\begin{tabular}{rcccc}
\toprule
\multicolumn{5}{c}{Comparação dos viés dos Estimadores $\hat{p_{1}}$ e $\tilde{p_{1}}$}\\
$n$ & $p$ & $B(\hat{p_{1}})$ & $B(\tilde{p_{1}})$ & $|B(\hat{p_{1}})|>|B(\tilde{p_{1}})|$ \\
\midrule
50 & 0.1 & 0.01158347 & 0.01108683 & TRUE \\
100 & 0.1 & 0.01147125 & 0.01124116 & TRUE \\
150 & 0.1 & 0.01129354 & 0.01113929 & TRUE \\
100000 & 0.1 & 0.01110255 & 0.01110307 & FALSE \\
\midrule
50 & 0.2 & 0.05069611 & 0.04952346 & TRUE \\
100 & 0.2 & 0.05113692 & 0.05050091 & TRUE \\
150 & 0.2 & 0.05021439 & 0.04981072 & TRUE\\
100000 & 0.2 & 0.05000023 & 0.05000045 & FALSE \\
\midrule
50 & 0.3 & 0.13303390 & 0.13062580 & TRUE \\
100 & 0.3 & 0.12839159 & 0.12709425 & TRUE \\
150 & 0.3 & 0.12910398 & 0.12821978 & TRUE \\
100000 & 0.3 & 0.12858837 & 0.12858547 & TRUE \\
\bottomrule
\end{tabular}
\end{table}
\end{frame}

\begin{frame}
\frametitle{Distribuição $BN(p,k)$}

\tiny
\begin{table}[h]
\caption{EQM dos estimadores $\hat{p_{1}}$ e $\tilde{p_{1}}$.}
\label{tab:p1EQM}
\centering
\begin{tabular}{rcccc}
\toprule
\multicolumn{5}{c}{Comparação dos EQM do Estimador $\hat{p_{1}}$ e $\tilde{p_{1}}$}\\
$n$ & $p$ & $EQM(\hat{p_{1}})$ & $EQM(\tilde{p_{1}})$ & $EQM(\hat{p_{1}})>EQM(\tilde{p_{1}})$ \\
\midrule
50 & 0.1 & 0.0001897657 & 0.0001792158 & TRUE \\
100 & 0.1 & 0.0001601663 & 0.0001552416 & TRUE \\
150 & 0.1 & 0.0001461537 & 0.0001429903 & TRUE \\
100000 & 0.1 & 0.0001232947 & 0.0001233068 & FALSE \\
\midrule
50 & 0.2 & 0.0028992563 & 0.0027791147 & TRUE \\
100 & 0.2 & 0.0027753925 & 0.0027122546 & TRUE \\
150 & 0.2 & 0.0026332882 & 0.0025933471 & TRUE\\
100000 & 0.2 & 0.0025001601 & 0.0025001831 & FALSE \\
\midrule
50 & 0.3 & 0.0188113709 & 0.0181701887 & TRUE \\
100 & 0.3 & 0.0170132255 & 0.0166791609 & TRUE \\
150 & 0.3 & 0.0170230361 & 0.0167979328 & TRUE \\
100000 & 0.3 & 0.0165355188 & 0.0165347806 & TRUE \\
\bottomrule
\end{tabular}
\end{table}
\end{frame}


\begin{frame}
\frametitle{Distribuição $BN(p,k)$}
\tiny
\begin{table}[h]
\caption{Comparação dos estimadores $\hat{p_{2}}$ e $\tilde{p_{2}}$.}
\label{tab:p2}
\centering
\begin{tabular}{rccc}
\toprule
\multicolumn{4}{c}{Comparação dos Estimadores $\hat{p_{2}}$ e $\tilde{p_{2}}$}\\
$n$ & $p$ & $|B(\hat{p_{2}})|>|B(\tilde{p_{2}})|$ & $EQM(\hat{p_{2}})>EQM(\tilde{p_{2}})$ \\
\midrule
50 & 0.1 & TRUE & FALSE \\
100 & 0.1 & TRUE & FALSE \\
150 & 0.1 & TRUE & FALSE \\
100000 & 0.1 & TRUE & FALSE \\
\midrule
50 & 0.2 & TRUE & TRUE \\
100 & 0.2 & TRUE & FALSE \\
150 & 0.2 & FALSE & FALSE\\
100000 & 0.2 & TRUE & FALSE \\
\midrule
50 & 0.3 & TRUE & FALSE \\
100 & 0.3 & TRUE & FALSE \\
150 & 0.3 & FALSE & FALSE \\
100000 & 0.3 & FALSE & FALSE \\
\bottomrule
\end{tabular}
\end{table}
\end{frame}

\begin{frame}
\frametitle{Distribuição $BN(p,k)$}
\tiny
\begin{table}[h]
\caption{Comparação dos estimadores $\hat{p_{3}}$ e $\tilde{p_{3}}$.}
\label{tab:p2}
\centering
\begin{tabular}{rccc}
\toprule
\multicolumn{4}{c}{Comparação dos Estimadores $\hat{p_{3}}$ e $\tilde{p_{3}}$}\\
$n$ & $p$ & $|B(\hat{p_{3}})|>|B(\tilde{p_{3}})|$ & $EQM(\hat{p_{3}})>EQM(\tilde{p_{3}})$ \\
\midrule
50 & 0.1 & TRUE & TRUE \\
100 & 0.1 & TRUE & TRUE \\
150 & 0.1 & TRUE & TRUE \\
100000 & 0.1 & TRUE & TRUE \\
\midrule
50 & 0.2 & TRUE & TRUE \\
100 & 0.2 & TRUE & TRUE \\
150 & 0.2 & TRUE & TRUE \\
100000 & 0.2 & TRUE & TRUE \\
\midrule
50 & 0.3 & TRUE & TRUE \\
100 & 0.3 & TRUE & TRUE \\
150 & 0.3 & TRUE & TRUE \\
100000 & 0.3 & TRUE & TRUE \\
\bottomrule
\end{tabular}
\end{table}
\end{frame}
\subsection{Distribuição Exponencial}
\begin{frame}
\frametitle{Distribuição Exponencial} 
A primeira tabela refere-se ao viés do estimador por máximo verossimilhança - $1/\bar{x}$.

\begin{table}[H]
\caption{Viés dos estimadores $\hat{\lambda_{1}}$ e $\tilde{\lambda_{1}}$.}
\label{tab:p1Vies}
\tiny
\centering
\begin{tabular}{rcccc}
\toprule
\multicolumn{5}{c}{Comparação dos viés do Estimador $\hat{\lambda_{1}}$}\\
$n$ & $\lambda$ & $B(\hat{\lambda_{1}})$ & $B(\tilde{\lambda_{1}})$ & $B(\hat{\lambda_{1}})>B(\tilde{\lambda_{1}})$ \\
\midrule
50 & 0.3 & 8.563419e-03 & 2.540733e-03  & TRUE \\
100 & 0.3 & 3.249035e-03 & 2.168678e-04 & TRUE \\
150 & 0.3 & 2.728954e-03 & 7.411608e-04 & TRUE \\
100000 & 0.3 & 1.972002e-05 & 1.913857e-05 & TRUE \\
\midrule
50 & 0.5 & 1.407817e-02 & 4.055892e-03  & TRUE \\
100 & 0.5 &  2.494619e-03 & -2.560141e-03 & TRUE \\
150 & 0.5 &  4.027130e-03  & 8.587662e-04 & TRUE\\
100000 & 0.5 & -3.786119e-05 & -4.269957e-05 & TRUE \\
\midrule
50 & 2 & 4.727426e-02 &  5.441983e-03 & TRUE \\
100 & 2 & 2.435143e-02 & 4.153628e-03 & TRUE \\
150 & 2 & 2.496612e-02 & 1.171509e-02 & TRUE \\
100000 & 2 & -5.638649e-05 & -5.874376e-05 & FALSE \\
\bottomrule
\end{tabular}
\end{table}
\end{frame}

\begin{frame}
\frametitle{Distribuição Exponencial}
A tabela a seguir trata-se dos valores do EQM em relação a máximo verossimilhança - $1/\hat{x}$.

\begin{table}[H]
\caption{EQM dos estimadores $\hat{\lambda_{1}}$ e $\tilde{\lambda_{1}}$.}
\label{tab:p1EQM}
\tiny
\centering
\begin{tabular}{rcccc}
\toprule
\multicolumn{5}{c}{Comparação dos EQM do Estimador $\hat{\lambda_{1}}$}\\
$n$ & $\lambda$ & $EQM(\hat{\lambda_{1}})$ & $EQM(\tilde{\lambda_{1}})$ & $EQM(\hat{\lambda_{1}})>EQM(\tilde{\lambda_{1}})$ \\
\midrule
50 & 0.3 & 2.194808e-03 & 2.056801e-03  & TRUE \\
100 & 0.3 & 8.528248e-04 & 8.278781e-04 & TRUE \\
150 & 0.3 & 5.894819e-04 & 5.826339e-04 & TRUE \\
100000 & 0.3 & 8.415198e-07 & 8.443515e-07 & FALSE \\
\midrule
50 & 0.5 & 6.094502e-03 & 5.732249e-03  & TRUE \\
100 & 0.5 &  2.729694e-03 & 2.662651e-03 & TRUE \\
150 & 0.5 &  1.750662e-03  & 1.727247e-03 & TRUE\\
100000 & 0.5 & 2.577277e-06 & 2.615424e-06 & FALSE \\
\midrule
50 & 2 & 8.901878e-02 &  8.467163e-02 & TRUE \\
100 & 2 & 4.237920e-02 &  4.133330e-02 & TRUE \\
150 & 2 & 3.086910e-02 & 3.009429e-02 & TRUE \\
100000 & 2 & 3.996599e-05 & 4.038770e-05 & FALSE \\
\bottomrule
\end{tabular}
\end{table}
\end{frame}

\begin{frame}
\frametitle{Distribuição Exponencial}
A tabela abaixo apresenta a diferença entre o viés no momento amostra de ordem 1.  $\sqrt{1/var(x)}$.


\begin{table}[H]
\caption{Viés dos estimadores $\hat{\lambda_{2}}$ e $\tilde{\lambda_{2}}$.}
\label{tab:p1Vies}
\tiny
\centering
\begin{tabular}{rcccc}
\toprule
\multicolumn{5}{c}{Comparação dos viés do Estimador $\hat{\lambda_{2}}$}\\
$n$ & $\lambda$ & $B(\hat{\lambda_{2}})$ & $B(\tilde{\lambda_{2}})$ & $B(\hat{\lambda_{2}})>B(\tilde{\lambda_{2}})$ \\
\midrule
50 & 0.3 & 2.000133e-02 & 3.046579e-03  & TRUE \\
100 & 0.3 & 8.539564e-03 & -3.095717e-04 & TRUE \\
150 & 0.3 & 5.592596e-03 & -4.209541e-04 & TRUE \\
100000 & 0.3 & 4.246102e-05 & 3.181654e-05 & TRUE \\
\midrule
50 & 0.5 & 3.125010e-02 & 3.198028e-03  & TRUE \\
100 & 0.5 &  1.661312e-02 & 2.302152e-03 & TRUE \\
150 & 0.5 &  1.293862e-02  & 2.823622e-03 & TRUE\\
100000 & 0.5 & 1.875328e-05 & -4.204228e-06 & TRUE \\
\midrule
50 & 2 & 1.312796e-01 &  2.056332e-02 & TRUE \\
100 & 2 & 4.087539e-02 & -1.814831e-02 & TRUE \\
150 & 2 & 5.083462e-02 & 9.218771e-03 & TRUE \\
100000 & 2 & 3.786189e-04 & 2.994606e-04 & FALSE \\
\bottomrule
\end{tabular}
\end{table}
\end{frame}

\begin{frame}
\frametitle{Distribuição Exponencial}
Tabela referente ao EQM do momento amostra de ordem 1.

\begin{table}[H]
\caption{EQM dos estimadores $\hat{\lambda_{2}}$ e $\tilde{\lambda_{2}}$.}
\label{tab:p1EQM}
\tiny
\centering
\begin{tabular}{rcccc}
\toprule
\multicolumn{5}{c}{Comparação dos EQM do Estimador $\hat{\lambda_{2}}$}\\
$n$ & $\lambda$ & $EQM(\hat{\lambda_{2}})$ & $EQM(\tilde{\lambda_{2}})$ & $EQM(\hat{\lambda_{2}})>EQM(\tilde{\lambda_{2}})$ \\
\midrule
50 & 0.3 & 4.174063e-03 & 3.976735e-03  & TRUE \\
100 & 0.3 & 1.990066e-03 & 1.991264e-03 & FALSE \\
150 & 0.3 & 1.247754e-03 & 1.253155e-03 & FALSE \\
100000 & 0.3 & 1.783617e-06 & 1.808211e-06 & FALSE \\
\midrule
50 & 0.5 & 1.201819e-02 & 1.144396e-02  & TRUE \\
100 & 0.5 &  5.240721e-03 & 5.232843e-03 & TRUE \\
150 & 0.5 &  3.648914e-03  & 3.671078e-03 & FALSE\\
100000 & 0.5 & 5.287958e-06 & 5.308073e-06 & FALSE \\
\midrule
50 & 2 & 1.903238e-01 &  1.831975e-01 & TRUE \\
100 & 2 & 8.064018e-02 &  8.440508e-02 & FALSE \\
150 & 2 & 5.738582e-02 & 5.732075e-02 & TRUE \\
100000 & 2 & 9.050068e-05 & 9.245068e-05 & FALSE \\
\bottomrule
\end{tabular}
\end{table}
\end{frame}

\begin{frame}
\frametitle{Distribuição Exponencial}
Quadro de comparação do viés do momento amostral de oredem 2. $\sqrt{1/mean^2}$

\begin{table}[H]
\caption{Viés dos estimadores $\hat{\lambda_{3}}$ e $\tilde{\lambda_{3}}$.}
\label{tab:p1Vies}
\tiny
\centering
\begin{tabular}{rcccc}
\toprule
\multicolumn{5}{c}{Comparação dos viés do Estimador $\hat{\lambda_{3}}$}\\
$n$ & $\lambda$ & $B(\hat{\lambda_{3}})$ & $B(\tilde{\lambda_{3}})$ & $B(\hat{\lambda_{3}})>B(\tilde{\lambda_{3}})$ \\
\midrule
50 & 0.3 & -8.009634e-02 & -8.679385e-02  & TRUE \\
100 & 0.3 & -8.220733e-02 & -8.578620e-02 & TRUE \\
150 & 0.3 & -8.501436e-02 & -8.756115e-02 & TRUE \\
100000 & 0.3 & -8.784743e-02 & -8.785105e-02 & TRUE \\
\midrule
50 & 0.5 & -1.360553e-01 & -1.467689e-01  & TRUE \\
100 & 0.5 &  -1.360536e-01 & -1.420749e-01 & TRUE \\
150 & 0.5 &  -1.424886e-01  & -1.466395e-01 & TRUE\\
100000 & 0.5 & -1.464497e-01 & -1.464554e-01 & TRUE \\
\midrule
50 & 2 & -5.376391e-01 &  -5.817923e-01 & TRUE \\
100 & 2 & -5.609116e-01 & -5.853201e-01 & TRUE \\
150 & 2 & -5.694220e-01 & -5.854173e-01 & TRUE \\
100000 & 2 & -5.856539e-01 & -5.856595e-01 & TRUE \\
\bottomrule
\end{tabular}
\end{table}
\end{frame}

\begin{frame}
\frametitle{Distribuição Exponencial}
Quadro com a comparação do EQM do momento amostral de oredem 2.

\begin{table}[H]
\caption{EQM dos estimadores $\hat{\lambda_{3}}$ e $\tilde{\lambda_{3}}$.}
\label{tab:p1EQM}
\tiny
\centering
\begin{tabular}{rcccc}
\toprule
\multicolumn{5}{c}{Comparação dos EQM do Estimador $\hat{\lambda_{3}}$}\\
$n$ & $\lambda$ & $EQM(\hat{\lambda_{3}})$ & $EQM(\tilde{\lambda_{3}})$ & $EQM(\hat{\lambda_{3}})>EQM(\tilde{\lambda_{3}})$ \\
\midrule
50 & 0.3 & 7.650300e-03 & 8.777607e-03  & FALSE \\
100 & 0.3 & 7.388421e-03 & 8.009593e-03 & TRUE \\
150 & 0.3 & 7.594362e-03 & 8.044213e-03 & TRUE \\
100000 & 0.3 & 7.717755e-03 & 7.718391e-03 & FALSE \\
\midrule
50 & 0.5 & 2.170059e-02 & 2.482344e-02  & FALSE \\
100 & 0.5 &  2.018802e-02 & 2.192635e-02 & FALSE \\
150 & 0.5 &  2.138781e-02  & 2.260983e-02 & FALSE\\
100000 & 0.5 & 2.144894e-02 & 2.145065e-02 & FALSE \\
\midrule
50 & 2 & 3.442212e-01 &  3.937709e-01 & FALSE \\
100 & 2 & 3.400614e-01 &  3.686495e-01 & FALSE \\
150 & 2 & 3.420241e-01 & 3.608734e-01 & FALSE \\
100000 & 2 & 3.430148e-01 & 3.430220e-01 & FALSE \\
\bottomrule
\end{tabular}
\end{table}
\end{frame}
%================================================================================

%=====================
\section{Conclusões}
%=====================

%================================================================================
\begin{frame}
\frametitle{Conclusões}

\end{frame}
%===============================================================================

\section{Referências Bibliográficas}

\begin{frame}
\frametitle{Referências Bibliográficas}


\bibliographystyle{plain}

\bibliography{Larangeiras}


\end{frame}
\end{document}